%==============================================================
%  lecons/analyse/gamma.tex — La fonction Gamma d'Euler
%==============================================================

\lecontitle{La fonction Gamma}{Analyse réelle — Intégrales à paramètre et fonctions spéciales}

%=== NOTATION LOCALE =========================================
\newcommand{\Gam}{\Gamma}
\newcommand{\intzi}{\int_0^{+\infty}}

%=============================================================
%  § 1 — DÉFINITION ET RÉSULTATS FONDAMENTAUX
%=============================================================
\secnum{navyblue}{1}{Définition et résultats fondamentaux}

\begin{tcolorbox}[thmbox]
  \textbf{Définition.}\enspace\index{fonction Gamma}
  Pour tout réel $x>0$, l'\textit{intégrale d'Euler de seconde espèce}
  \[
    \Gam(x) \;=\; \intzi t^{x-1}\,e^{-t}\,\mathrm{d}t
  \]
  est convergente. Elle définit la \textit{fonction Gamma} sur $]0,+\infty[$.

  \medskip
  \textbf{Théorème (propriétés fondamentales).}
  \begin{itemize}
    \item \textit{Équation fonctionnelle :} pour tout $x>0$,
    \[
      \Gam(x+1) = x\,\Gam(x), \qquad \Gam(1)=1,
    \]
    d'où $\Gam(n+1)=n!$ pour tout entier $n\geq 0$ : $\Gam$ prolonge la
    factorielle.
    \item \textit{Régularité :} $\Gam$ est de classe $\mathcal{C}^\infty$ sur
    $]0,+\infty[$ (dérivation sous le signe intégral), avec
    \[
      \Gam^{(k)}(x) = \intzi {(\ln t)}^k\,t^{x-1}\,e^{-t}\,\mathrm{d}t.
    \]
    \item \textit{Log-convexité :} $\ln\circ\,\Gam$ est convexe sur
    $]0,+\infty[$. Le \textit{théorème de Bohr--Mollerup}
    \index{théorème de Bohr--Mollerup} caractérise $\Gam$ comme l'unique
    fonction $f:\,]0,+\infty[\,\to\,]0,+\infty[$ vérifiant $f(1)=1$,
    $f(x+1)=x\,f(x)$ et $\ln f$ convexe.
    \item \textit{Valeurs remarquables :}
    \[
      \Gam\!\left(\tfrac12\right)=\sqrt{\pi}, \qquad
      \Gam(x)\,\Gam(1-x)=\frac{\pi}{\sin(\pi x)} \quad (0<x<1)
    \]
    (\textit{formule des compléments}\index{formule des compléments}).
    \item \textit{Équivalent de Stirling :}
    \[
      \Gam(x+1) \;\sim\; \sqrt{2\pi x}\,{\left(\frac{x}{e}\right)}^{\!x}
      \quad (x\to+\infty).
    \]
  \end{itemize}
\end{tcolorbox}

\begin{significationintuitive}
La factorielle $n!$ n'est définie que pour des entiers ; pourtant l'identité
$\,n! = n\cdot(n-1)!$ et l'intégrale $\intzi t^{n}e^{-t}\,\mathrm{d}t=n!$
suggèrent un sens « continu ». La fonction $\Gam$ réalise ce prolongement :
elle interpole les points $(n,\,(n-1)!)$ par une courbe lisse, et l'exigence
de log-convexité (théorème de Bohr--Mollerup) en fait le prolongement
\emph{canonique}, à l'exclusion de toute autre interpolation raisonnable.
\end{significationintuitive}

\decsep{}

%=============================================================
%  § 2 — DÉMONSTRATION
%=============================================================
\secnum{navyblue}{2}{Démonstration}

\begin{ideepreuve}
La convergence se traite séparément en $0$ (où $t^{x-1}$ peut être singulier)
et en $+\infty$ (où $e^{-t}$ domine). L'équation fonctionnelle vient d'une
intégration par parties. La régularité $\mathcal{C}^\infty$ repose sur le
théorème de dérivation sous le signe intégral, justifié par une domination
locale uniforme. Enfin $\Gam(1/2)$ se ramène à l'intégrale de Gauss via le
changement $t=u^2$.
\end{ideepreuve}

\vspace{10pt}
\rubriq{Preuve}

\begin{mdframed}[style=proofbar]
  \textit{Étape 1 — Convergence.}
  Soit $x>0$. La fonction $t\mapsto t^{x-1}e^{-t}$ est continue positive sur
  $]0,+\infty[$. Au voisinage de $0$, $t^{x-1}e^{-t}\sim t^{x-1}$, intégrable
  ssi $x-1>-1$, c'est-à-dire $x>0$. Au voisinage de $+\infty$, on a
  $t^{x-1}e^{-t}=o(t^{-2})$ (croissances comparées), donc intégrable. Ainsi
  $\Gam(x)$ est bien défini pour tout $x>0$, et l'intégrale diverge pour
  $x\leq 0$.

  \medskip
  \textit{Étape 2 — Équation fonctionnelle.}
  Pour $x>0$, une intégration par parties sur $[\varepsilon,A]$ avec
  $u=t^{x}$, $v'=e^{-t}$ donne
  \[
    \int_\varepsilon^A t^{x}e^{-t}\,\mathrm{d}t
    = {\big[-t^{x}e^{-t}\big]}_\varepsilon^A
      + x\int_\varepsilon^A t^{x-1}e^{-t}\,\mathrm{d}t.
  \]
  Le terme tout intégré ${\big[-t^{x}e^{-t}\big]}_\varepsilon^A
  =\varepsilon^{x}e^{-\varepsilon}-A^{x}e^{-A}$ tend vers $0$ quand
  $\varepsilon\to 0^+$ et $A\to+\infty$ : en $0$ car $\varepsilon^{x}\to 0$
  (puisque $x>0$), en $+\infty$ par croissances comparées. En passant à la
  limite :
  \[
    \Gam(x+1) = x\,\Gam(x).
  \]
  De plus $\Gam(1)=\intzi e^{-t}\,\mathrm{d}t=1$, d'où par récurrence
  $\Gam(n+1)=n\,\Gam(n)=\dots=n!\,\Gam(1)=n!$.

  \medskip
  \textit{Étape 3 — Continuité et dérivabilité $\mathcal{C}^\infty$.}
  Posons $f(x,t)=t^{x-1}e^{-t}$. Fixons un segment $[a,b]\subset\,]0,+\infty[$.
  La dérivée partielle $\partial_x^{\,k} f(x,t)={(\ln t)}^k\,t^{x-1}e^{-t}$ est
  continue en $x$, et pour tout $x\in[a,b]$ on a la domination
  \[
    \bigl|{(\ln t)}^k\,t^{x-1}\bigr|\,e^{-t}
    \;\leq\; {|\ln t|}^k\bigl(t^{a-1}+t^{b-1}\bigr)\,e^{-t}
    \;=:\;\varphi_k(t),
  \]
  où $\varphi_k$ est intégrable sur $]0,+\infty[$ (le facteur logarithmique ne
  modifie ni la convergence en $0$ — car $t^{\delta}{|\ln t|}^k\to 0$ pour tout
  $\delta>0$ — ni celle en $+\infty$, dominée par $e^{-t}$). Le théorème de
  dérivation sous le signe intégral s'applique sur $[a,b]$, donc $\Gam$ est
  $\mathcal{C}^\infty$ et
  \[
    \Gam^{(k)}(x) = \intzi {(\ln t)}^k\,t^{x-1}e^{-t}\,\mathrm{d}t.
  \]
  En particulier $\Gam''(x)=\intzi {(\ln t)}^2\,t^{x-1}e^{-t}\,\mathrm{d}t>0$,
  donc $\Gam$ est convexe ; la log-convexité s'obtient par l'inégalité de
  Cauchy--Schwarz appliquée à
  $\Gam'(x)^2\leq\Gam(x)\,\Gam''(x)$, qui équivaut à
  ${(\ln\Gam)}''\geq 0$.

  \medskip
  \textit{Étape 4 — Valeur $\Gam(1/2)$ et intégrale de Gauss.}
  Le changement de variable $t=u^2$ (donc $\mathrm{d}t=2u\,\mathrm{d}u$) donne
  \[
    \Gam\!\left(\tfrac12\right)
    = \intzi t^{-1/2}e^{-t}\,\mathrm{d}t
    = \intzi u^{-1}e^{-u^2}\,2u\,\mathrm{d}u
    = 2\intzi e^{-u^2}\,\mathrm{d}u
    = \sqrt{\pi},
  \]
  où l'on a utilisé l'intégrale de Gauss $\intzi e^{-u^2}\,\mathrm{d}u
  =\tfrac{\sqrt{\pi}}{2}$ — elle-même limite des intégrales de Wallis
  $W_n=\int_0^{\pi/2}\sin^n t\,\mathrm{d}t$, puisque
  $\int_0^{\sqrt{n}}{(1-u^2/n)}^n\,\mathrm{d}u=\sqrt{n}\,W_{2n+1}\to
  \tfrac{\sqrt{\pi}}{2}$ d'après l'équivalent $W_n\sim\sqrt{\pi/(2n)}$ établi
  dans la leçon sur les intégrales de Wallis. On en
  déduit, via l'équation fonctionnelle,
  $\Gam(3/2)=\tfrac12\sqrt{\pi}$, $\Gam(5/2)=\tfrac34\sqrt{\pi}$, etc.

  \smallskip
  La formule des compléments et l'équivalent de Stirling, de preuve
  sensiblement plus délicate, sont admis dans cette leçon.
  \hfill$\blacksquare$
\end{mdframed}

\decsep{}

%=============================================================
%  § 3 — EXEMPLES, CONTRE-EXEMPLES ET EXERCICES
%=============================================================
\secnum{exgreen}{3}{Exemples, contre-exemples et exercices}

\rubriq{Exemples}

\begin{mdframed}[style=exbar]
  \textbf{1. Valeurs aux entiers et demi-entiers.}
  \[
    \Gam(n+1)=n!,\qquad
    \Gam\!\left(\tfrac12\right)=\sqrt{\pi},\qquad
    \Gam\!\left(\tfrac32\right)=\tfrac{\sqrt{\pi}}{2},\qquad
    \Gam\!\left(\tfrac52\right)=\tfrac{3\sqrt{\pi}}{4}.
  \]
  Plus généralement
  $\Gam\!\left(n+\tfrac12\right)=\dfrac{(2n)!}{4^n\,n!}\sqrt{\pi}
   =\dfrac{(2n-1)!!}{2^{n}}\sqrt{\pi}$.

  \smallskip\noindent
  \textbf{2. Fonction Bêta.}\enspace\index{fonction Bêta}
  Pour $x,y>0$,
  \[
    B(x,y)=\int_0^1 u^{x-1}{(1-u)}^{y-1}\,\mathrm{d}u
          =\frac{\Gam(x)\,\Gam(y)}{\Gam(x+y)}.
  \]
  C'est l'outil de calcul d'innombrables intégrales : par exemple
  $\displaystyle\int_0^1\frac{\mathrm{d}u}{\sqrt{u(1-u)}}=B\!\left(\tfrac12,\tfrac12\right)
   =\dfrac{\Gam(1/2)^2}{\Gam(1)}=\pi$.

  \smallskip\noindent
  \textbf{3. Lien avec les intégrales de Wallis.}\enspace\index{intégrales de Wallis}
  Le changement $u=\sin^2 t$ donne
  $W_n=\int_0^{\pi/2} \sin^n t\,\mathrm{d}t=\tfrac12 B\!\left(\tfrac{n+1}{2},\tfrac12\right)
   =\dfrac{\sqrt{\pi}}{2}\,\dfrac{\Gam\!\left(\frac{n+1}{2}\right)}
   {\Gam\!\left(\frac{n}{2}+1\right)}$.

  \smallskip\noindent
  \textbf{4. Stirling factoriel.}\enspace\index{formule de Stirling}
  En $x=n$, l'équivalent de Stirling redonne
  $n!=\Gam(n+1)\sim\sqrt{2\pi n}\,{(n/e)}^n$.
\end{mdframed}

\vspace{6pt}
\rubriq{Contre-exemples et pièges}

\begin{mdframed}[style=cexbar]
  \textbf{$\Gam$ n'est pas définie en $0$ ni aux entiers négatifs.}\enspace%
  L'intégrale $\intzi t^{x-1}e^{-t}\,\mathrm{d}t$ \emph{diverge} dès que
  $x\leq 0$ (divergence en $0$). Ainsi $\Gam(0)$ « n'existe pas » au sens de
  l'intégrale.

  \smallskip\noindent
  \textbf{Pôles du prolongement.}\enspace%
  La relation $\Gam(x)=\dfrac{\Gam(x+n)}{x(x+1)\cdots(x+n-1)}$ permet de
  prolonger $\Gam$ aux réels (puis complexes) non entiers négatifs : on obtient
  une fonction \textit{méromorphe} dont les seuls pôles sont $0,-1,-2,\dots$,
  de résidu $\dfrac{{(-1)}^n}{n!}$ en $-n$. Il ne faut donc \emph{pas} croire
  que $\Gam$ est partout finie.

  \smallskip\noindent
  \textbf{Convexité seule ne suffit pas.}\enspace%
  Une fonction interpolant la factorielle et seulement \emph{convexe} n'est pas
  unique : c'est la log-convexité (Bohr--Mollerup) qui force l'unicité.
\end{mdframed}

\vspace{6pt}
\exolabel

\begin{mdframed}[style=exobar]
  \begin{enumerate}
    \item Calculer $\displaystyle\intzi t^{3}e^{-t}\,\mathrm{d}t$ et
          $\displaystyle\intzi \sqrt{t}\,e^{-t}\,\mathrm{d}t$.
          \textit{Indication :} reconnaître $\Gam(4)$ et $\Gam(3/2)$.

    \item Montrer que pour $a>0$,
          $\displaystyle\intzi t^{x-1}e^{-at}\,\mathrm{d}t
          =\dfrac{\Gam(x)}{a^{x}}$ (changement de variable affine).

    \item Établir $\displaystyle\intzi e^{-t^2}\,t^{2n}\,\mathrm{d}t
          =\tfrac12\,\Gam\!\left(n+\tfrac12\right)
          =\dfrac{(2n-1)!!}{2^{n+1}}\sqrt{\pi}$.

    \item À partir de la formule des compléments, déduire la valeur de
          $\displaystyle\intzi\frac{\mathrm{d}u}{1+u^4}$.
          \textit{Indication :} via $u^4=v$, se ramener à
          $\tfrac14 B\!\left(\tfrac14,\tfrac34\right)$ puis
          $\Gam(1/4)\Gam(3/4)=\pi/\sin(\pi/4)$.

    \item (Bohr--Mollerup) Soit $f:\,]0,+\infty[\,\to\,]0,+\infty[$ vérifiant
          $f(1)=1$, $f(x+1)=x f(x)$ et $\ln f$ convexe. Montrer que
          $f=\Gam$. \textit{Indication :} encadrer $\ln f(x)$ pour
          $0<x<1$ à l'aide de la convexité de $\ln f$ entre des entiers, puis
          faire tendre $n\to+\infty$.

    \item (Formule de Gauss) Établir
          $\displaystyle\Gam(x)=\lim_{n\to+\infty}
          \frac{n!\,n^{x}}{x(x+1)\cdots(x+n)}$ pour $x>0$, et en déduire la
          formule du produit
          $\dfrac{1}{\Gam(x)}=x\,e^{\gamma x}\prod_{k\geq 1}
          \left(1+\tfrac{x}{k}\right)e^{-x/k}$
          (avec $\gamma$ la constante d'Euler).
  \end{enumerate}
\end{mdframed}

\leconfooter{L.\ Euler (1729) ; H.\ Bohr \& J.\ Mollerup (1922)}
